\begin{abstract}
% Differential privacy (DP) and federated learning (FL) are combined as advanced privacy-preserving methods when training on-device language models in production mobile keyboard applications. DP-Follow-the-Regularized-Leader (DP-FTRL) algorithms, leveraging correlated noise mechanisms such as tree aggregation or matrix factorization, are widely used in practice for their superior privacy-utility trade-off and compatibility with FL systems. This paper presents a novel variant of DP-FTRL by adapting the recent theoretical advancements of the Buffered Linear Toeplitz (BLT) mechanism to multi-participant scenarios. In the FL setting, our BLT mechanism demonstrates enhanced privacy-utility trade-off and improved memory efficiency than the widely used tree aggregation mechanism. Moreover, BLT achieves comparable privacy and utility to the state-of-the-art banded matrix factorization mechanism, while significantly simplifying usage requirements and reducing memory. The flexibility of the BLT mechanism allows seamless integration with existing DP FL implementations in production environments. We evaluate BLT-DP-FTRL on the StackOverflow dataset, serving as a research simulation benchmark, and across four on-device language model tasks in a production FL system. Our empirical results highlight the advantages of the BLT mechanism to elevate the practicality and effectiveness of DP in real-world scenarios. 

The state-of-the-art for training on-device language models for mobile keyboard applications combines federated learning (FL) with differential privacy (DP) via the DP-Follow-the-Regularized-Leader (DP-FTRL) algorithm. Two variants of DP-FTRL are used in practice, tree aggregation and matrix factorization. However, tree aggregation suffers from significantly suboptimal privacy/utility tradeoffs, while matrix mechanisms require expensive optimization parameterized by hard-to-estimate-in-advance constants, and high runtime memory costs.
This paper extends the recently introduced Buffered Linear Toeplitz (BLT) mechanism to multi-participation scenarios. Our BLT-DP-FTRL maintains the ease-of-use advantages of tree aggregation, while essentially matching matrix factorization in terms of utility and privacy. 
We evaluate BLT-DP-FTRL on the StackOverflow dataset, serving as a re-producible simulation benchmark, and across four on-device language model tasks in a production FL system. Our empirical results highlight the advantages of the BLT mechanism and elevate the practicality and effectiveness of DP in real-world scenarios. 

\end{abstract}